% Generalization-check table: do the Construction Scene sweep's headline
% findings (analysis/lhs_sweep_findings.txt) hold on a different scene?
% Compares Construction's 150-point sweep against a smaller 50-point first
% look at PublicTransit Scene - see
% analysis/lhs_sweep_publictransit_findings.txt for the full writeup this
% table summarizes. Both Dialogue priority, Advanced (Summed-bus) ducking.
%
% Requires packages already in main.tex's preamble: booktabs (siunitx not
% needed here - see below). No \cellcolor/colortbl or longtable needed;
% this is a small, non-ranked comparison table, not a per-row data dump.
%
% Plain right-aligned (r) columns, not siunitx S columns: the two data
% columns mix very different value shapes (signed dB levels, unsigned
% probabilities, signed correlation coefficients) that don't share one
% table-format spec cleanly, and this table is small enough that manual
% alignment costs nothing - same reasoning as parameter_sweep_table.tex's
% choice to avoid S columns where they'd add fragility without real
% benefit.
%
% All figures recomputed directly from analysis/lhs_sweep_results.csv and
% analysis/lhs_sweep_results_publictransit_dialogue.csv for this table (not
% hand-copied from the two findings.txt files) to guarantee both columns
% were computed identically. The two "range across tested threshold" rows
% use each sweep's actual most-aggressive (~-60dB) and least-aggressive
% (~-18dB) tested point - not the single best-margin-gain point - so the
% range reflects threshold's real effect, not a cherry-picked outcome.
% Baseline margin/STOI: Construction from analysis/folder_export_results.csv
% (Construction/Dialogue row, already using the true-stem-sum baseline fix
% and the 0.85-octave margin padding); Public Transit from
% run_lhs_sweep.py's own baseline computation (same methodology, same
% script).

\begin{table}[h]
\centering
\footnotesize
\caption{Generalization check: key parameter-sweep results, Construction
Scene (150-point sweep) vs. Public Transit Scene (50-point sweep), both
Dialogue priority, Advanced (Summed-bus) ducking. Public Transit's sweep
is smaller -- a first look, not yet sized to cite per-tier settings the
way Construction's Aggressive/Balanced/Conservative tiers are -- but the
results already land in the same range.}
\label{tab:scene-generalization}
\begin{tabular}{@{}l r r@{}}
\toprule
 & \textbf{Construction} & \textbf{Public Transit} \\
\midrule
Sweep size (points) & 150 & 50 \\
Unprocessed baseline margin (dB) & -9.51 & -13.28 \\
Unprocessed baseline STOI & 0.393 & 0.275 \\
Margin gain, most aggressive threshold tested\textsuperscript{a} (dB) & +6.79 & +6.28 \\
Margin gain, least aggressive threshold tested\textsuperscript{a} (dB) & +0.05 & +0.03 \\
STOI gain, most aggressive threshold tested\textsuperscript{a} & +0.106 & +0.094 \\
STOI gain, least aggressive threshold tested\textsuperscript{a} & +0.003 & +0.002 \\
Loudness reduction, most aggressive threshold\textsuperscript{a} (pp) & +11.1 & +12.7 \\
Loudness reduction, least aggressive threshold\textsuperscript{a} (pp) & +0.1 & +0.0 \\
\bottomrule
\end{tabular}
\par\vspace{4pt}
\raggedright\footnotesize\textsuperscript{a} Most aggressive $\approx$
\SI{-60}{\decibel}, least aggressive $\approx$ \SI{-18}{\decibel} -- the
two ends of the search range, same in both sweeps. Both scenes show the
same pattern concretely: pushing threshold from the weak end to the
aggressive end moves margin gain from essentially \SI{0}{\decibel} to
roughly \SI{+6.5}{\decibel}, and STOI gain and loudness reduction move by
a comparable proportion right along with it in both scenes -- not just a
similar correlation statistic, but the same real improvement available at
the aggressive end, and the same near-total loss of effect at the weak
end. Public Transit's unprocessed baseline is noticeably worse (more
heavily masked before any processing) than Construction's, making this a
meaningfully different, not merely repeated, test case for the same
pattern to hold on.
\end{table}
